\section{Proofs for Appendix~\ref{sec:repack}}\label{app:repack}
Notation and conventions are those of Appendix~\ref{sec:repack}.

\subsection{The exchange calculus}
\begin{proof}[Proof of Lemma~\ref{lem:rpex}]
The pairs of $T$ are not soft pairs of $H-Q$, and the members of $H-Q$ have pairwise disjoint pair sets; hence $H'=(H-Q)+T$ is modification-disjoint, and $T\notin H-Q$ (a member of $H-Q$ has soft pairs), so $|H'|=|H|$ (if $T=Q$ there is nothing to prove). As graphs on the node set $H-Q$, $I(H')-T=I(H)-Q$, an induced subgraph of $I(H)$. Adding one node to a graph raises the treewidth by at most one (put it in every bag), so $\tw(I(H'))\le\tw(I(H)-Q)+1\le\tw(I(H))+1$. For $j$ exchanges the same argument applies with $H-\{Q_1,\dots,Q_j\}$, the $T_i$ being pairwise pair-disjoint and disjoint from its soft pairs, and $j$ added nodes.

Now let $(G,H,\ell)$ be a YES instance. Since $|H'|=|H|$, $\OPT\le|H'|+\ell$, so $(G,H',\ell)$ is a YES instance. Let $M$ be a maximal family hard over $H'$. By Lemma~\ref{lem:fam}, $|M|\le\ell$, by Lemma~\ref{lem:MAX} $H''=H'\cup M$ is pair-maximal, $\tw(I(H''))\le\tw(I(H'))+|M|\le\tw(I(H))+1+|M|$, and $(G,H'',\ell-|M|)$ is equivalent to $(G,H',\ell)$, hence to $(G,H,\ell)$. That $H'$ need not be pair-maximal is shown by the computer test of Appendix~\ref{sec:repack}.
\end{proof}

\begin{proof}[Proof of Lemma~\ref{lem:rptwin}]
Suppose $a\in C_1$, $b\in C_2$ with $C_1\ne C_2$. Since $N[a]=N[b]$, $ab\in E$, and $ab$ is an edge across clusters, so it is edited. Let $c_a$ be the number of pairs $\{a,x\}$, $x\notin\{a,b\}$, that disagree with $S$ (that is, are an edge across clusters or a non-edge inside a cluster), and $c_b$ likewise. These pairs together with $ab$ contribute $c_a+c_b+1$ to the cost; all other pairs do not involve $a$ or $b$. Move $a$ into $C_2$. For $x\notin\{a,b\}$, $a$ and $x$ have the same adjacency as $b$ and $x$, and $a$ is now in the cluster of $b$, so $\{a,x\}$ disagrees exactly when $\{b,x\}$ did before: $c_b$ pairs; the pairs at $b$ are unchanged, $c_b$, and $ab$ is now an edge inside a cluster and agrees. The new cost of these pairs is $2c_b$, all other pairs are unchanged. Symmetrically, moving $b$ into $C_1$ gives $2c_a$. Since $\min\{2c_a,2c_b\}\le c_a+c_b<c_a+c_b+1$, $S$ is not optimal.
\end{proof}

\begin{proof}[Proof of Corollary~\ref{cor:rpP3}]
Let $ab\in S$ be inserted: $a,b$ lie in one cluster $C$ and $ab\notin E$. By Lemma~\ref{lem:half}(b) there is $z\in C$ with $az,zb\in E$, so $(a,z,b)$ is an induced $P_3$ containing $ab$. Let $ab\in S$ be deleted: $ab\in E$ and $a,b$ lie in different clusters. By Lemma~\ref{lem:rptwin}, $a,b$ are not closed twins, so some $x\notin\{a,b\}$ is adjacent to exactly one of them, say to $a$; then $x\,a\,b$ is an induced $P_3$ containing $ab$ (adjacent to $b$ instead: $x\,b\,a$).
\end{proof}

\begin{proof}[Proof of Lemma~\ref{lem:rpwit}]
By Corollary~\ref{cor:rpP3} witnesses exist. Let $T$ be a witness; the pair $e$ is hard, so $T$ has at most two soft pairs. As $H$ is pair-maximal, $T$ does not have all three pairs hard, so it has at least one soft pair. If $e$ is inserted, a non-edge of $G$, then in the induced $P_3$ $T$ the pair $e$ joins the two ends and the opposite vertex $z$ is the centre. If $e$ is deleted, an edge of $G$, then $e$ is an edge of $T$ and $z$ is an end.

(a) Let $T$ have two soft pairs. They are the two pairs other than $e$, namely $\{a,z\}$ and $\{b,z\}$; let $P,P'$ own them, so $a,z\in V(P)$ and $b,z\in V(P')$. If $P=P'$ then $a,b\in V(P)$ and $e\in\pairs(P)$ would be soft, a contradiction. So $P\ne P'$, $z\in V(P)\cap V(P')$, and $z$ is a common soft neighbour of $a$ and $b$ ($z\ne a,b$): $e$ is bridged.

(b) Otherwise every witness has exactly one soft pair $s$, owned by $Q$, and its other two pairs (one of them $e$) are hard. So the pairs of $T$ are not soft pairs of $H-Q$, and Lemma~\ref{lem:rpex} shows that $(H-Q)+T$ is a packing of size $|H|$; $e\in\pairs(T)$ is soft in it.
\end{proof}

\begin{proof}[Proof of Corollary~\ref{cor:rpabs}]
If $e$ is not bridged, no witness has two soft pairs by Lemma~\ref{lem:rpwit}(a) (whose conclusion needs both ends covered), so case (b) applies to any witness, and one exists. The width bound is Lemma~\ref{lem:rpex}.
\end{proof}

\begin{proof}[Proof of Lemma~\ref{lem:rpunc}]
Let $e=\{r,r'\}\in U$ with $r,r'$ uncovered. By Corollary~\ref{cor:rpP3} it has a witness $T=\{r,r',x\}$. The pairs $\{r,x\}$ and $\{r',x\}$ are hard, since $r$ and $r'$ lie in no member, and $e$ is hard. So $T$ is an induced $P_3$ with all pairs hard, contradicting pair-maximality.
\end{proof}

\begin{proof}[Proof of Lemma~\ref{lem:rpdisp}]
The pair $az$ joins two vertices of the cluster $C$ of $a,b$ and is an edge of $G$ (a step of the witness), so it is not in $S$. $Q$ is an induced $P_3$ of $G$ on $\{a,z,r\}$ with the edge $az$, so exactly one of $ar,zr$ is an edge of $G$ and the other a non-edge. If $r\in C$, all three pairs of $Q$ are edges of $G\sd S$: the edge among $ar,zr$ is kept, and the non-edge is inserted. If $r\notin C$, then $ar,zr$ join different clusters: the edge is deleted and the non-edge is untouched. In both cases $S\cap\pairs(Q)=\{f\}$ with $f\in\{ar,zr\}$. Also $r\ne b$: otherwise $\pairs(Q)=\pairs(T)\ni e$ would be soft. The pairs of $T$ in $S$ are $e$ only ($az\notin S$, and $zb$ is an edge inside $C$). So relative to $H'=(H-Q)+T$ the hard edits are $(U\setminus\{e\})\cup\{f\}$: $e$ is soft, $f\notin\pairs(T)$ is hard. The ends of $f$ are $r$ and $a$ or $z$. As $r\notin V(T)$, $r$ is covered after the exchange iff it lies in a member of $H-Q$.
\end{proof}

\begin{proof}[Proof of Lemma~\ref{lem:rpout}]
All pairs at $r$ are hard, so $N_h(r)=N(r)$. Let $z\in N(r)$ and $w\in N_h(z)$, $w\ne r$. Then $rz$ and $zw$ are hard edges, so by Lemma~\ref{lem:HE}(a) (centre $z$) $rw\in E$ or $rw$ is soft; $rw$ is not soft as $r$ is uncovered, so $w\in N(r)$. Hence the hard component of $r$ is $N[r]$. For $z,z'\in N(r)$, Lemma~\ref{lem:HE}(a) with centre $r$ gives: $zz'\in E$ (a hard edge or a soft pair) or $zz'$ is soft. A soft pair lies in one member, so its ends are co-bagged; a pair of $N[r]$ that is not co-bagged is therefore not soft, and is an edge, hence a hard edge.
\end{proof}

\subsection{Proofs of Propositions~\ref{prop:rpFAR} and~\ref{prop:rpSp}}
\begin{proof}[Proof of Proposition~\ref{prop:rpFAR}]
(a) An induced $P_3$ has exactly two edges. The members of a packing have disjoint pair sets, hence are edge-disjoint. As $|E(G_m)|=2m+1$, a packing has at most $m$ members, and $H_m$ has $m$. $\OPT=m+1$ is Proposition~\ref{prop:FAR}(ii). By Lemma~\ref{lem:decomp} every packing $H^*$ has $\ell^*=|H|+\ell-|H^*|\ge\OPT-\nu_{\mathrm{mod}}=1$ in a YES instance of excess $1$.

(b) Pairs of $P_i$: $\{2i,2i+1\},\{2i+1,2i+2\},\{2i,2i+2\}$. $T=(1,2m-1,2m-2)$: $1\sim2m-1$ (chord), $2m-1\sim2m-2$ (path edge), and $1\not\sim2m-2$ because the neighbours of $1$ are $0,2,2m-1$ and $2m-2\ge4$ for $m\ge3$. So $T$ is an induced $P_3$ with centre $2m-1$ and pairs $\{1,2m-1\}$, $\{2m-2,2m-1\}$, $\{1,2m-2\}$. The first is in no $P_i$ ($2m-2\ge4$ apart, both odd); the second lies only in $P_{m-1}$; the third is neither a consecutive pair nor an even-even pair ($2m-3\ge3$). So $T$ avoids the pairs of $P_0,\dots,P_{m-2}$ and Lemma~\ref{lem:rpex} applies with $Q=P_{m-1}$. Soft edges of $H^*$: the path edges $\{j,j+1\}$, $j\le2m-3$, and $\{2m-2,2m-1\}$, and the chord. The only hard edge is $\{2m-1,2m\}$; an induced $P_3$ has two edges, at most one hard, so it contains a soft pair: $H^*$ is pair-maximal. $V(T)=\{1,2m-2,2m-1\}$ meets $V(P_i)=\{2i,2i+1,2i+2\}$, $i\le m-2$, exactly for $i=0$ (vertex $1$) and $i=m-2$ (vertex $2m-2$); and $V(P_i)\cap V(P_{i+1})=\{2i+2\}$. So $I(H^*)$ is the cycle $P_0,P_1,\dots,P_{m-2},T,P_0$ on $m\ge3$ nodes, of treewidth $2$.

(c) The partition has the edits: the path edges $\{2i+1,2i+2\}$ for $0\le i\le m-1$ and the chord, $m+1$ pairs, and no inserted pair; so its cost is $m+1=\OPT$. Relative to $H^*$ the edits $\{2i+1,2i+2\}$, $i\le m-2$, lie in $P_i$, the chord lies in $T$; the only hard edit is $\{2m-1,2m\}$ and $2m\notin V(H^*)=\{0,\dots,2m-1\}$. The vertex $2m$ has the single neighbour $2m-1$, whose neighbours are $2m-2,2m,1$. So $\mathrm{Ball}(2m)\cap V(H^*)\subseteq\{2m-1,2m-2,1\}=V(T)$, pairwise co-bagged in every decomposition (one member), and $2m\in L_B$ (hard degree $1$). So $2m$ is a local hub and $\{2m-1,2m\}$ is eligible; this optimum lies in $\mathcal S$ for $H^*$, $\ell^*=1$.
\end{proof}

\begin{proof}[Proof of Proposition~\ref{prop:rpSp}]
We use the description of $\mathrm{Sp}(F)$ in Proposition~\ref{prop:five}. \emph{$Y$ is an induced $P_3$:} $yu_0,u_0b_0\in E$ ($b_0,u_0\in\mathcal K_*^+$) and $yb_0\notin E$ ($y$ is adjacent only to $u_0,c_0$). Its pairs $\{y,u_0\}$, $\{y,b_0\}$ are hard ($y$ is uncovered) and $\{u_0,b_0\}\in\pairs(X_0)$ only. \emph{$Z$ is an induced $P_3$:} $b_0u_1$ is a spine edge, $u_1c_1$ is an edge inside $\mathcal K^+_{h(1)}$ (both belong to it), and $b_0c_1\notin E$, because $c_1$ has no edge outside $\mathcal K^+_{h(1)}$ and $b_0\in\mathcal K^+_*$ with $*\ne h(1)$ (only $e_0$ has head $*$). The soft edges are $a_ib_i,b_ic_i,u_ib_i,b_iu_{i+1}$ and the soft non-edges are $a_ic_i$, $u_iu_{i+1}$; so $\{u_1,c_1\}$ and $\{b_0,c_1\}$ are hard, and $\{b_0,u_1\}\in\pairs(X_0)$ only. The pairs of $Y$ and $Z$ are disjoint, and none is a soft pair of $H-X_0$. So $H'=(H-X_0)+Y+Z$ is a packing and $|H'|=|H|+1$. By Lemma~\ref{lem:decomp} $\OPT\ge|H'|=|H|+1$, and $\OPT\le|H|+1$ by the partition of Proposition~\ref{prop:five}(S2).

\emph{Width.} $I(H)$ is a caterpillar with spine $X_0,\dots,X_N$ and leaf $P_i$ at $X_i$. Removing $X_0$ isolates $P_0$. $Y$ meets only $P_0$ and $Z$ (all through $b_0$; $u_0$ lies in $X_0$ and in no other member, and $y$ in no member of $H$). $Z$ meets $P_0,Y$ (via $b_0$), $X_1$ (via $u_1$) and $P_1$ (via $c_1$). So $I(H')$ is obtained from the caterpillar $X_1,\dots,X_N$ with leaves $P_1,\dots,P_N$ by adding the triangle $\{P_0,Y,Z\}$ and joining $Z$ to $X_1$ and $P_1$ (the triangle $\{Z,X_1,P_1\}$, since $X_1P_1$ is an edge). Its blocks are triangles and edges, so $\tw(I(H'))\le2$.

\emph{Class.} The optimum of (S2) has the single hard edit $\{y,b_0\}$ with $y$ uncovered. All hard degrees are at most $7\le B=11$ (S3), so $y,b_0\in L_B$. $\mathrm{Ball}(y)=\{y,u_0,c_0,b_0\}$ (the neighbours of $u_0$ and $c_0$ are $\mathcal K_*^+$ and $y$). The covered vertices $u_0,b_0\in V(X_0)$ and $c_0,b_0\in V(P_0)$, and $X_0,P_0$ are adjacent through $b_0$, so they are co-bagged in every decomposition: $y$ is a local hub and the edit is eligible.
\end{proof}

\subsection{Proof of Theorem~\ref{thm:rpPL} and Corollary~\ref{cor:rpstatic}}
Fix $H$, $\ell$ and a decomposition; \emph{eligibility} (Definition~\ref{def:rpclass}) is then computable from $G,H$ and the decomposition. For a covered $v\in L_B$ let $\mathrm{Elig}(v)$ be the set of partners $p$ with $\{v,p\}$ eligible; $|\mathrm{Elig}(v)|\le|\mathrm{Ball}(v)|\le Z^2$. Eligibility of a covered pair is symmetric. Every eligible pair has a covered end in $L_B$.

\emph{Labels and status.} A \emph{labelling} assigns to every member $P$ a state $x_P$ among the five of Lemma~\ref{lem:five}, and to every covered $v\in L_B$ a set $E_v\subseteq\mathrm{Elig}(v)$ with $|E_v|\le\ell$. It defines $S=\bigcup_Px_P\cup U$ with $U=\{\{v,p\}:p\in E_v\}$ and is \emph{symmetric} if for covered $v,v'$ with $\{v,v'\}$ eligible, $v'\in E_v\iff v\in E_{v'}$. The status of a pair in $G\sd S$ is $G$-status xor membership in $S$. At a bag $t$, the status of a pair is \emph{determined by the state at $t$} if it is soft with owner in $\beta_t$ (status from $x_P$), or hard and not eligible (status of $G$), or hard and eligible with an end $p\in X_t\cap L_B$ (status from $[q\in E_p]$).

\begin{lemma}[Visibility]\label{lem:rpvis}
Let $H$ be pair-maximal. For every vertex triple $\Delta$ of $G$ there is a bag $t$ at which the statuses of all three pairs of $\Delta$ are determined by the state, or $\Delta$ has no soft pair and no eligible pair (then all its statuses are static and $\Delta$ is not an induced $P_3$ of $G$).
\end{lemma}
\begin{proof}
\emph{$\Delta$ has a soft pair.} The owners of the soft pairs of $\Delta$ form a clique of $I(H)$ (Observation~\ref{obs:two}), so some bag $t$ contains all of them (Helly). Every vertex of $\Delta$ incident to a soft pair lies in $X_t$. A hard pair of $\Delta$ has an end among these vertices: if $\Delta$ has one soft pair $\{a,b\}$, the hard pairs are $\{a,c\},\{b,c\}$; if it has two or three, all vertices of $\Delta$ lie in $X_t$. So each hard pair is static or has an end in $X_t$ with a label, and is determined at $t$ (a vertex outside $L_B$ is in no eligible pair).

\emph{$\Delta$ has no soft pair and no eligible pair.} Then $\Delta$ has three hard pairs, all static, and it is not an induced $P_3$ of $G$ by pair-maximality.

\emph{$\Delta=\{p,q,c\}$ has no soft pair and the eligible pair $\{p,q\}$.} (i) $p,q$ covered, co-bagged: take $t\in\mathcal T_p\cap\mathcal T_q$; the pairs $\{p,c\},\{q,c\}$ are static or determined by the labels of $p$ and $q$. (ii) $p=r$ uncovered, $q=v$ covered, $r$ a local hub. The pair $\{r,v\}$ is determined by the label of $v$. If $c$ is uncovered then $\{r,c\}$ is not eligible, hence static, and $\{v,c\}$ is static or determined by the label of $v$: any bag of $\mathcal T_v$ works. If $c$ is covered and $c\in\mathrm{Ball}(r)$, then $v,c\in V(H)\cap\mathrm{Ball}(r)$ are co-bagged; take $t\in\mathcal T_v\cap\mathcal T_c$: $\{r,c\}$ is determined by the label of $c$, $\{v,c\}$ by that of $v$. If $c$ is covered and $c\notin\mathrm{Ball}(r)$, then $\{r,c\}$ is static and a bag of $\mathcal T_v$ works. By Lemma~\ref{lem:rpunc} no eligible pair has two uncovered ends (by definition of eligibility an eligible pair has a covered end), so the cases are exhaustive.
\end{proof}

\begin{proof}[Proof of Theorem~\ref{thm:rpPL}]
\emph{Algorithm.} Take a nice tree decomposition of $I(H)$ of width $w$. A state at a bag $t$ is $(x_P)_{P\in\beta_t}\in[5]^{\beta_t}$ together with $E_v\subseteq\mathrm{Elig}(v)$, $|E_v|\le\ell$, for every $v\in X_t\cap V(H)\cap L_B$. A vertex $v$ is introduced when the first member through $v$ enters, its label is chosen there, it is carried unchanged through the subtree $\mathcal T_v$, and forgotten with the last member through $v$; a join node requires equal states of the two children. At each bag the table keeps, for each state, the minimum cost, where the cost is $\sum_P|x_P|+\sum_v|E_v\setminus V(H)|+\frac12\sum_v|E_v\cap V(H)|$, each term charged at the introduction of its member or vertex. A state is \emph{consistent at $t$} if (S) for $v,v'\in X_t\cap L_B$ covered: $v'\in E_v\iff v\in E_{v'}$; and (V) for every triple $\Delta$ whose three statuses are determined at $t$, $\Delta$ is not an induced $P_3$ of $G\sd S$, where the statuses are those given by the state. Inconsistent states are discarded. The output $m^*$ is the minimum cost at the root, over all states in the root table (and $\infty$ if none).

\emph{Number of states.} $|\beta_t|\le w+1$ gives $5^{w+1}$ choices of $(x_P)$; $|X_t|\le3(w+1)$ vertices have a label, each a set of at most $\ell$ elements of a set of at most $Z^2$ elements, at most $(Z^2+1)^\ell$ choices. Together $5^{w+1}(Z^2+1)^{3\ell(w+1)}$ states. Each transition and each test (enumerating triples) takes $n^{O(1)}$ time, and $Z=3D+3\ell+1\le6(w+1)+3\ell+1$ since $D=2c\le2(w+1)$ (Lemma~\ref{lem:deg}).

\emph{Soundness.} A consistent labelling of the whole decomposition gives labels $E_v$ constant along $\mathcal T_v$ (carried through, equal at joins) and states $x_P$; define $S=\bigcup x_P\cup U$ with $U=\{\{v,p\}:p\in E_v\}$. Two covered vertices $v,v'$ with $\{v,v'\}$ eligible are co-bagged, so some bag contains both, and (S) there gives $v'\in E_v\iff v\in E_{v'}$ globally: a covered pair of $U$ is seen from both ends, so the cost term $\frac12\sum_v|E_v\cap V(H)|$ counts it once, and $|S|$ equals the cost. By Lemma~\ref{lem:rpvis}, every triple is either static and not an induced $P_3$ of $G$ (so unaffected by $S$) or has all three statuses determined at some bag, where it is verified. Statuses at different bags agree, since they come from the same global labelling (the symmetry of labels makes the status of an eligible covered pair independent of the end that is used). Hence no triple is an induced $P_3$ of $G\sd S$: $S$ is a solution, all of whose hard edits are eligible, each covered vertex in at most $\ell$ of them: $S\in\mathcal S$, and $m^*$ is at least $\min\{|S|:S\in\mathcal S\}\ge\OPT$.

\emph{Completeness.} Let $S\in\mathcal S$. Put $x_P=S\cap\pairs(P)$, which is one of the five states because $V(P)$ is a cluster graph in $G\sd S$ (Lemma~\ref{lem:decomp}, Lemma~\ref{lem:five}), and $E_v=\{p:\{v,p\}\in S\text{ hard}\}\subseteq\mathrm{Elig}(v)$, $|E_v|\le\ell$. The labelling is symmetric by construction, its cost is $|S|$, and every triple is not an induced $P_3$ of $G\sd S$; so all states are consistent and $m^*\le|S|$. Hence $m^*=\min\{|S|:S\in\mathcal S\}$. If some optimal solution of size at most $|H|+\ell$ is in $\mathcal S$ then $m^*=\OPT$. The last sentence of the theorem follows.
\end{proof}
The restriction of the labels to $\mathrm{Elig}(v)$ is used twice: it makes (S) checkable in a bag common to both ends, and it gives the local-hub condition that makes the triples at an uncovered vertex visible; an uncovered $r$ lies in no bag, so two edits at $r$ with partners that are not co-bagged could not be compared.

\begin{proof}[Proof of Corollary~\ref{cor:rpstatic}]
Let a decomposition co-bag every pair of $\Lambda(H)$. Let $S$ be an optimal solution of size at most $|H|+\ell$ and $\{x,y\}\in U$ with $x$ covered. By Corollary~\ref{cor:cand}, $x,y\in L_B$ and $y\in\mathrm{Ball}(x)$; so if $y$ is covered, $\{x,y\}\in\Lambda(H)$ (take $r=x$), hence co-bagged and eligible. If $y$ is uncovered, $y\in L_B$ is a hub: any two covered $a,b\in\mathrm{Ball}(y)$ lie in $\Lambda(H)$ (take $r=y$), so they are co-bagged, and $\{y,x\}$ is eligible. Lemma~\ref{lem:rpunc} excludes two uncovered ends, and $|U|\le\ell$. So $S\in\mathcal S$ and Theorem~\ref{thm:rpPL} is exact. The equivalence with the treewidth of $I(H)$ plus links: a decomposition of $I(H)$ that puts a member $P\ni a$ and a member $P'\ni b$ in a common bag is a decomposition of $I(H)$ plus the link $PP'$, and conversely a decomposition of $I(H)$ plus links is one of $I(H)$ in which the linked members share a bag.
\end{proof}

\subsection{Proofs for the limits of repacking}
\begin{proof}[Proof of Theorem~\ref{thm:rpN}]
(a) At $\ell=0$ CEaMP is NP-hard~\cite[Thm.~1]{LPS23}, even with every vertex in at most $23$ members (Theorem~\ref{thm:eth}). Suppose the algorithm $\mathcal R$ exists with constants $s=s(0)$, $f=f(0)$. Given $(G,H,0)$, run $\mathcal R$ and let $H^*$ be its output. The guarantee holds only on YES instances, so validate: (i) test that $H^*$ is a modification-disjoint packing of induced $P_3$'s of $G$, else reject; (ii) compute, in time $f^{O(f^3)}|H^*|$~\cite{Bod96}, a tree decomposition of $I(H^*)$ of width at most $f$, else reject; (iii) put $\ell^*=|H|-|H^*|$ and reject unless $0\le\ell^*\le s$. Finally run the CSP algorithm (Theorem~\ref{thm:A}) on $(G,H^*,\ell^*)$ with this decomposition, in time $5^{f+1}n^{2s+O(1)}$, polynomial, and answer accordingly. Correctness: for a valid packing $H^*$, $(G,H^*,\ell^*)$ is a YES instance iff $\OPT\le|H^*|+\ell^*=|H|$, i.e.\ iff $(G,H,0)$ is. So whenever a decision is made it is correct; the rejections in (i)--(iii) are unreachable on YES instances (where the guarantees hold, and $\ell^*\ge0$ because $|H^*|\le\OPT\le|H|$), and are correct on NO instances. Hence NP is solved in polynomial time. The theorem used is the algorithm of Section~\ref{sec:csp}, valid for every packing; the majority-closure theorem (Theorem~\ref{thm:sfA}) is not used, and the soft triangles in the gadgets of~\cite{LPS23} play no role.

(b) Suppose $\mathcal R$ maps every YES instance $(G,H,0)$ to $H^*$ with $|H^*|=|H|$ and $\tw(I(H^*))\le|H|^{1-\varepsilon}$. Take the instances of the proof of Theorem~\ref{thm:eth} from a sparsified 3-SAT formula on $N$ variables, with $|H|=O(N)$. Run $\mathcal R$, validate that $H^*$ is a packing of induced $P_3$'s with $|H^*|=|H|$ (else reject), compute a tree decomposition of $I(H^*)$ of width at most $2\tw+1$ in time $2^{O(|H|^{1-\varepsilon})}|H|$~\cite{Kor21}, rejecting if its width exceeds $2|H|^{1-\varepsilon}+1$, and apply Theorem~\ref{thm:l0} with it, in time $3^{2|H|^{1-\varepsilon}+3}n^{O(1)}$. Since $(G,H^*,0)$ is equivalent to $(G,H,0)$ for a valid packing of the same size, this decides 3-SAT in time $2^{O(N^{1-\varepsilon})}$ (plus the $2^{o(N)}$ of sparsification and the polynomial running time of $\mathcal R$), contradicting ETH.

The two remarks follow: a tight instance ($\ell=0$, $\nu_{\mathrm{mod}}=\OPT=|H|$) has no packing larger than $H$, so ``large treewidth implies a better packing'' fails on it, and tight instances of unbounded treewidth exist under ETH by Theorem~\ref{thm:eth} (they are the instances of its proof).
\end{proof}

\begin{proof}[Proof of Proposition~\ref{prop:rpnf}]
The induced $P_3$'s of $G$ are $(x,h,y)$ with $x,y$ leaves and $h$ a hub (a leaf and both hubs form a triangle). Every member has a hub as its centre. Two members with the same centre $h$ and leaf pairs $\{x,y\},\{x',y'\}$ have disjoint pair sets only if $\{x,y\}\cap\{x',y'\}=\emptyset$, impossible among three leaves. So each hub is the centre of at most one member, and $\nu_{\mathrm{mod}}\le2$. A packing of size $2$ has one member with centre $3$ and leaf pair $\{x,y\}$ and one with centre $4$ and a different leaf pair $\{x',y'\}$; two different $2$-subsets of three leaves cover all leaves, so every maximum packing covers all five vertices ($\nu_{\mathrm{mod}}=2$, six maximum packings). Enumerating all $52$ partitions gives $\OPT=3$ and the four optima $\{0,1,2,3,4\}$, $\{0\},\{1,2,3,4\}$, $\{1\},\{0,2,3,4\}$, $\{2\},\{0,1,3,4\}$; for each of the six maximum packings each optimum has a hard edit (all vertices are covered, so it is covered-covered). For example, for $H^*=\{(0,3,1),(1,4,2)\}$ and the optimum $\{0\},\{1,2,3,4\}$ the hard edit is $\{0,4\}$ and $0,4$ have the common soft neighbour $1$: it is bridged. The counts quoted in Section~\ref{sec:repack} come from an exhaustive enumeration of all connected graphs with at least one induced $P_3$, over all maximum packings and all optima.
\end{proof}

\begin{proof}[Proof of Proposition~\ref{prop:rpaug}]
The hard edges of $H_0$ are $\{0,1\}$ and $\{2m+1,2m+2\}$, which do not share a vertex; an induced $P_3$ of the path is three consecutive vertices, with two edges, at most one hard, so it contains a soft pair: $H_0$ is pair-maximal. A packing has at most $\lfloor(2m+2)/2\rfloor=m+1$ members (two edges each, edge-disjoint), and $\{(2i,2i+1,2i+2):0\le i\le m\}$ attains it; $\OPT\le m+1$ by the partition $\{0,1\},\dots,\{2m,2m+1\},\{2m+2\}$, so $\nu_{\mathrm{mod}}=\OPT=m+1$ and $\ell=1$ for $H_0$ ($|H_0|=m$), $w=1$. A packing of size $m+1$ must use all $2m+2$ edges. The edge $\{0,1\}$ lies only in $(0,1,2)$, so this is a member; then $\{2,3\}$ lies only in $(2,3,4)$ among the remaining ones, and so on: the packing is unique, $\{(2i,2i+1,2i+2)\}$, and it shares no member with $H_0$ (whose members start at odd vertices). An improving exchange removes $r$ members of $H_0$ and adds $r+1$ members pair-disjoint from the rest; the result is a packing of size $m+1$, hence the unique one, so all $m$ members of $H_0$ are removed: $r=m$.
\end{proof}

\begin{proof}[Proof of Proposition~\ref{prop:rpLOC}]
Run the algorithm of Open Problem~\ref{op:LOC}, validate that $H^*$ is a pair-maximal packing of induced $P_3$'s with $\ell^*\ge0$ (else reject, correct because then $\OPT\ge|H^*|>|H|+\ell$ or the output is invalid on a NO instance), and run Theorem~\ref{thm:rpPL} on $(G,H^*,\ell^*)$ with the decomposition of width at most $g(\ell,w)$; the time is $5^{g+1}(Z^2+1)^{3\ell^*(g+1)}n^{O(1)}$, a function of $(\ell,w)$ times a polynomial. Accept iff $m^*\le|H^*|+\ell^*$. Soundness holds for every output by Theorem~\ref{thm:rpPL}, and on a YES instance some optimum lies in the class, so $m^*=\OPT\le|H^*|+\ell^*=|H|+\ell$.
\end{proof}
